\scanpage{177}
\mspage{35}

est $W$-asphérique. En fait, $Y_{/f_!(e_{X^\wedge})}$ est \inserted{\unsure{elle}} $=$ une catégorie filtrante.

En effet, le foncteur $f_!$ commute aux limites inductives, \deleted{donc} $f_!e_{X^\wedge}$ est limite inductive filtrante des $f_!(x_i)=f(x_i)$. Donc $Y_{/f_!(e_{X^\wedge})}$ est limite inductive filtrante des $Y_{/f(x_i)}$, Donc \deleted{il est} \inserted{c'est} $=$ une catégorie filtrante, a fortiori \deleted{\unsure{elle est}} $W$-asphérique.

Donc \unsure{finalement|facilement}, on trouve des exemples non triviaux de catégories tot. $W$-asphériques, \unsure{i.e. des} exemples dans des catégories sans objet final.

% slanted
\reversemarginpar\margin{déjà dans SGA4 I 8 \rule{0.4pt}{4ex}}\normalmarginpar

\uline{Corollaire} 2.18. Soit $X$ dans \unsure{Cat}, $F$ dans $X^\wedge$. Pour que $F$ soit ind-représentable, il f. et s. que $X_{/F}$ soit filtrante.

En effet, \inserted{par (2.15;)} $X_{/F}$ est filtré ssi l'objet final $F$ de $X_{/F}$ \deleted{est limite inductive} \inserted{provient d'un} \deleted{d'une famille filtrante} ind-objet de $X_{/F}$, mais cela signifie aussi que $F$ est la limite inductive dans $X^\wedge$ d'un ind-objet de $X$, i.e. que $F$ est ind-représentable.

% archive sheet number: 177
